\chapter{Main Counting Formulas (Prop 10.11, 10.12)}
\label{chap:counting}

These are the \textbf{main unconditional counting theorems} of [BMSZb]
Section~10. They express the PBP count via a recursion over the dual
partition. The statements only require shape / dual-partition legality; no
auxiliary proof objects appear.

\section{Recursive counting functions}

\begin{definition}[Tail class for D type]
  \label{def:tailClass_D}
  \lean{tailClass_D, TailClass}
  \leanok
  \uses{def:PBP}
  The \emph{tail class} of $\tau \in \PBP_D$ is one of $\{\mathrm{DD}, \mathrm{RC}, \mathrm{SS}\}$
  determined by the tail symbol (Def.~\ref{def:tail}): $x_\tau = d$ gives $\mathrm{DD}$,
  $x_\tau \in \{r, c\}$ gives $\mathrm{RC}$, $x_\tau = s$ gives $\mathrm{SS}$.
\end{definition}

\begin{definition}[countPBP\_D — recursion on D-type dp]
  \label{def:countPBP_D}
  \lean{countPBP_D, countPBP_D_of_YD, countPBP_D_of_YD_aux, tripleSum}
  \leanok
  \uses{def:DualPart, def:tailClass_D, prop:10_9_D, prop:10_8_C_prim}
  $\mathrm{countPBP\_D}(dp) \in \mathbb{N}^3$ is a triple
  $(n_{\mathrm{DD}}, n_{\mathrm{RC}}, n_{\mathrm{SS}})$ counting PBPs by tail class.

  Base case: $dp = []$ gives $(0, 0, 0)$ (or $(1, 0, 0)$ for the empty orbit, depending on
  the initial convention; see the Lean definition).

  Step: remove the first part of $dp$, recursively count on the remaining orbit, and
  combine via descent relations (Prop~\ref{prop:10_8_C_prim}, Prop~\ref{prop:10_8_C_bal},
  or the balanced D split).

  $\mathrm{tripleSum}(n_{\mathrm{DD}}, n_{\mathrm{RC}}, n_{\mathrm{SS}}) := n_{\mathrm{DD}} + n_{\mathrm{RC}} + n_{\mathrm{SS}}$ is the total count.
\end{definition}

\begin{definition}[countPBP\_B — recursion on B-type dp]
  \label{def:countPBP_B}
  \lean{countPBP_B}
  \leanok
  \uses{def:DualPart, def:countPBP_D, prop:10_8_M_prim, prop:10_8_M_bal}
  $\mathrm{countPBP\_B}(dp) \in \mathbb{N}^3$ via analogous recursion, using
  the M$\to$B descent (Prop~10.8 M) plus recursion to D on a shifted shape.
\end{definition}

\begin{definition}[countPBP\_C — recursion on C-type dp]
  \label{def:countPBP_C}
  \lean{countPBP_C}
  \leanok
  \uses{def:countPBP_D, prop:10_8_C_prim, prop:10_8_C_bal}
  $\mathrm{countPBP\_C}(dp) \in \mathbb{N}$ via descent to D:
  \[
    \mathrm{countPBP\_C}(dp) =
    \begin{cases}
      \text{(primitive formula)} & \mu_Q.\mathrm{colLen}(0) \geq \mu_P.\mathrm{colLen}(0) \\
      n_{\mathrm{DD}} + n_{\mathrm{RC}} & \text{balanced}
    \end{cases}
  \]
\end{definition}

\begin{definition}[countPBP\_M — recursion on M-type dp]
  \label{def:countPBP_M}
  \lean{countPBP_M}
  \leanok
  \uses{def:countPBP_B, prop:10_8_M_prim, prop:10_8_M_bal}
  $\mathrm{countPBP\_M}(dp) \in \mathbb{N}$ via descent to B.
\end{definition}

\section{Main counting theorems}

\begin{proposition}[Prop 10.11 B — counting for B type]
  \label{prop:10_11_B}
  \lean{card_PBPSet_B_eq_tripleSum_countPBP_B, prop_10_11_B, card_B_combined}
  \leanok
  \uses{cor:10_10, def:countPBP_B, prop:10_8_M_prim, prop:10_8_M_bal}
  For any valid B-type dual partition $dp$ with shape $(\mu_P, \mu_Q)$ derived from $dp$:
  \[ |\PBP_{B^+}(\mu_P, \mu_Q)| + |\PBP_{B^-}(\mu_P, \mu_Q)| = \mathrm{tripleSum}\big(\mathrm{countPBP\_B}(dp)\big). \]
\end{proposition}
\begin{proof}
  \leanok
  Joint induction on the ``Total'', ``A1'', ``A3'' components following the
  B-type lifting structure (see \texttt{CorrespondenceB.lean}, main theorem
  \texttt{card\_B\_combined}). The balanced cases use Prop~\ref{prop:10_8_M_bal}
  and a fiber decomposition of $\PBP_{B^\pm}$ along bottom-$\mathcal{Q}$ paint.
  Primitive cases use Prop~\ref{prop:10_8_M_prim} and the M-type balanced
  bijection together with induction hypothesis on the shifted $\mu_P$.
\end{proof}

\begin{proposition}[Prop 10.11 D — counting for D type]
  \label{prop:10_11_D}
  \lean{card_PBPSet_D_eq_tripleSum_countPBP_D, prop_10_11_D}
  \leanok
  \uses{cor:10_10, def:countPBP_D}
  For any valid D-type dual partition $dp$:
  \[ |\PBP_D(\mu_P, \mu_Q)| = \mathrm{tripleSum}\big(\mathrm{countPBP\_D}(dp)\big). \]
\end{proposition}
\begin{proof}
  \leanok
  See the Lean formalization referenced on the statement.
\end{proof}


\begin{proposition}[Prop 10.12 C — counting for C type]
  \label{prop:10_12_C}
  \lean{card_PBPSet_C_eq_countPBP_C, prop_10_12_C}
  \leanok
  \uses{prop:10_8_C_prim, prop:10_8_C_bal, prop:10_11_D, def:countPBP_C, prop:10_2_C}
  For any valid C-type dual partition $dp$ (with $dp \neq []$):
  \[ |\PBP_C(\mu_P, \mu_Q)| = \mathrm{countPBP\_C}(dp). \]
\end{proposition}
\begin{proof}
  \leanok
  See the Lean formalization referenced on the statement.
\end{proof}


\begin{proposition}[Prop 10.12 M — counting for M type]
  \label{prop:10_12_M}
  \lean{card_PBPSet_M_eq_countPBP_M, prop_10_12_M}
  \leanok
  \uses{prop:10_8_M_prim, prop:10_8_M_bal, prop:10_11_B, def:countPBP_M, prop:10_2_M}
  For any valid M-type dual partition $dp$:
  \[ |\PBP_M(\mu_P, \mu_Q)| = \mathrm{countPBP\_M}(dp). \]
\end{proposition}
\begin{proof}
  \leanok
  See the Lean formalization referenced on the statement.
\end{proof}


\section{Unconditional status}

The counting recursion formulas (\ref{prop:10_11_B}, \ref{prop:10_11_D},
\ref{prop:10_12_C}, \ref{prop:10_12_M}) are \textbf{fully unconditional in the
following sense}: the hypotheses are only shape/dual-partition legality
(\texttt{SortedGE}, parity, positivity, shape match), all of which are
defining properties of ``valid B/D/C/M dual partition''. No auxiliary proof
object (ACResult, ILS, L, twist equality) is needed.
